
\chapter{Example C Client}
\label{cha:logclient}

This chapter explains the usage of the C programming interface
with an example program called {\tt LogClient}. This program
is able to read the {\tt OB.log} file created by ConceptBase server
and performs the operations stored in this file.

The full source code of this program is in the directory
\verb+$CB_HOME/examples/Clients/LogClient+. The file
{\tt LogClient.c} contains beside the main program some little
functions to read the log file. This source file must be compiled
and linked together with a version of the ConceptBase library {\tt libCB.a}.
The file {\tt MakeLogClient} is a makefile, which executes the
necessary commands to compile and link the file with {\tt gcc} on a Unix-platform.
The following paragraphs explain only the important parts of the main program.

Before any functions of libCB can be used, one must include the header
file {\tt CBinterface.h}.
\begin{verbatim}
  #include <CBinterface.h>

  int main(int argc,char* argv[]) {

    int PortNr;
    char* HostName;
    char* UserName;
    Answer *ans;
    Server *gserver;
    char* command;
    char* arg;

    char   *ClientName  = "LogClient";
\end{verbatim}
The variables {\tt PortNr}, {\tt HostName}, {\tt UserName} and {\tt ClientName}
are initialized with the command line arguments and passed to the \verb+connect_CB_server+
function below. {\tt ans} stores the answer of an operation with the ConceptBase
server. {\tt gserver} is a pointer to a {\tt Server} structure which is filled by the
connect-function. {\tt command} is the command which has been read from the log file and
{\tt arg} is the argument for this command.

\begin{verbatim}
    /* Reading and checking command line arguments */
    /* ... */

    /* Connect to CBserver */
    connect_CB_server(PortNr, HostName, ClientName, UserName, &gserver);
    if (!gserver) {
        fprintf(stderr,"Connection failed!\n");
        return 1;
    }
\end{verbatim}
The function \verb+connect_CB_server+ opens an IPC socket to the specified
ConceptBase server and performs an {\tt ENROLL\_ME} method as described in
chapter \ref{serverinterface}.
If the connection can be successfully established, the {\tt gserver} variable
points to the connected server. Otherwise, {\tt gserver} will be {\tt NULL}.

Now, the program begins to read the log file. As long as there are
commands in the log file, the variable {\tt command} points to a string
containing the actual method and {\tt arg} contains the arguments of
this method.

Depending on the value of {\tt command}, the program executes
the corresponding function to pass the method with its arguments {\tt arg}
to the ConceptBase server. Possible values for {\tt command} are
e.g.\ {\tt tellCB, untell,ask\_frames, ...}

\begin{verbatim}
    /* Read commands from logfile until end of file */
    while(readLogCommand(fp,&command,&arg)) {

        /* Ask user, if the command should be executed */
        /* ... */

        /* Tell */
        if (!strcmp(command,"tell")) {
            printf("Telling: %s \n\n",arg);
            ans=tellCB( gserver, arg );
        }

        /* Untell */
        if (!strcmp(command,"untell")) {
            printf("Untelling: %s \n\n",arg);
            ans=untell( gserver, arg );
        }

        /* Tell Model */
        if (!strcmp(command,"tell_model")) {
            printf("Loading models: %s \n\n",arg);
            files=commaList2charArray(arg);
            ans=tell_model( gserver, files );
            for(i=0;i<MAX_FILES;i++) {
                if (files[i])
                    free(files[i]);
                free(files);
            }
        }
\end{verbatim}
Note, that the {\tt tell} and {\tt untell} functions take a simple string
containing frames as argument, whereas the function {\tt tell\_model}
takes a list of filenames as argument. The frames are loaded from these
files by ConceptBase and the told to the knowledge base.

Tell and untell operations return a pointer to an {\ttfamily Answer} object.
For tell and untell, it is sufficient to check the completion value of the answer.
The {\ttfamily return\_data} can be ignored for these methods.

The following {\tt ask} functions return also an {\ttfamily Answer} object.
The answer of the query is stored in the field {\ttfamily return\_data}, the
completion is {\tt CB\_OK}, if the query could be evaluated. Otherwise
the completion will {\ttfamily CB\_ERROR} or {\ttfamily CB\_TIMEOUT}.

\begin{verbatim}
        /* Ask objnames */
        if (!strcmp(command,"ask_objnames")) {
            printf("Ask (OBJNAMES): %s \n\n",arg);
            ans=ask_objnames( gserver, arg, "LABEL","Now" );
            printf("Answer: %s\n\n", ans->return_data);
        }

        /* Ask frames */
        if (!strcmp(command,"ask_frames")) {
            printf("Ask (OBJNAMES): %s \n\n",arg);
            ans=ask_frames( gserver, arg, "LABEL","Now" );
            printf("Answer: %s\n\n", ans->return_data);
        }

        /* Check completion */
        if (ans && ans->completion) {
            fprintf(stderr,
                    ">>> Server reports error on method: %s(%s)\n\n",
                    command,arg);
        }
\end{verbatim}
When an error occured, i.e. {\tt completion} is not zero (another value than {\tt CB\_OK}),
than a error message is printed on the console. Perhaps, it is also useful to get the
all error messages from ConceptBase server, but this is not done here\footnote{But that
should be done in a {\em good} client program.}.
\begin{verbatim}
    }
    /* Close connection to CBserver */
    disconnect_CB_server(gserver);

    return 0;
}
\end{verbatim}
If the while-loop is finished the connection to the ConceptBase server can be
closed with the function \verb+disconnect_CB_server+ and the program is finished.

